Der neue Abstand ist $r_2 = (1+p)\,r_1$ mit $p = \zuO$. Da die Gravitationskraft umgekehrt proportional zum Quadrat des Abstands ist, gilt
\begin{align}
 \frac{F_2}{F_1} &= \left(\frac{r_1}{r_2}\right)^2 = \fakF \\
   &= \frac{1}{\left(1+\zuO\right)^2} = \fak
\end{align}
Die relative Änderung der Kraft ist
\begin{align}
 \frac{F_2-F_1}{F_1} &= \frac{F_2}{F_1}-1 = \fak-1 \approx \result{\dFP{0}{2}}
\end{align}
Die Gravitationskraft nimmt prozentual knapp doppelt so stark ab, wie der Abstand zunimmt.
